\section{The model and its distinguished divisor}\label{sec:model}

We verify the local ring in Theorem~\ref{thm:main} and construct the
divisor on which the integral class will be lifted. The special fiber
forces a factor of five in residue degrees away from $m$, while the
divisor through $m$ has residue degree one and a uniformizer of norm
$\pi$. These properties will distinguish its class from classes
whose supporting points specialize away from $m$. For a finite
extension $E/L$, write $\OO_E$ for its valuation ring and $k_E$
for its residue field.

\begin{lemma}\label{lem:model}
The generic fiber $C=\X_L$ is an integral projective curve. The
ring $A=\OO_{\X,m}$ is a two-dimensional Noetherian regular local
domain of mixed characteristic $(0,5)$ with
$\pi\in\mathfrak m_A^2$. For every closed point $q\in C$
with residue field $E=\kappa(q)$, the $E$-point extends to
$\Spec\OO_E\to\X$. If its closed image is not $m$, then
$5\mid[k_E:\F_5]$.

On the chart $Z\ne0$, set $x=X/Z$, $y=Y/Z$, and let $t$ be the
image of $x$ in $S=A/yA$. Then
\begin{equation}\label{eq:divisor}
 S\cong V[t]/\bigl((t-1)^5+\zeta\bigr)
\end{equation}
is a DVR with uniformizer $t$. Its fraction field $D$ is totally
ramified of degree five over $L$, and
\begin{equation}\label{eq:norm-t}
 N_{D/L}(t)=\pi,\qquad \eta=1-t,\qquad \eta^5=\zeta.
\end{equation}
The equation $Y=0$ cuts out a single closed point $q_0\in C$
with residue field $D$.
\end{lemma}

\begin{proof}
The polynomial
\[
 \Phi_5(1-T)=T^4-5T^3+10T^2-10T+5
\]
is Eisenstein at five. Therefore $\pi=1-\zeta$ is a uniformizer
of $V$, its residue field is $\F_5$, and $5=u\pi^4$ for some
$u\in V^\times$.
The homogeneous equation $P$ defining $\X$ reduces to
\[
 \overline P=X^5-XY^4+Y^5.
\]
If $\alpha$ is a root of $s^5-s+1$, then Frobenius sends
$\alpha$ to $\alpha-1$. Its orbit has length exactly five:
the $r$th iterate is $\alpha-r$, which first returns at $r=5$.
Thus $s^5-s+1$ is irreducible over $\F_5$, and so is its
homogenization $\overline P$. This remains irreducible on adjoining
the independent variable $Z$. A factorization of the primitive
homogeneous polynomial $P$ over $L$ would, by Gauss' lemma, give
primitive positive-degree homogeneous factors over $V$ and hence
a factorization of $\overline P$. Consequently $C$ is integral.
It is a projective plane curve, so it has dimension one.

Put
\[
 F=(x-1)^5-xy^4+y^5+\zeta,\qquad
 R=V[x,y]/(F),\qquad \mathfrak m=(\pi,x,y)\subset R.
\]
In the regular local ring $V[x,y]_{(\pi,x,y)}$ of dimension
three, $F\equiv-\pi$ modulo the square of the maximal ideal.
Thus quotienting by $F$ gives the regular local ring
$A=R_{\mathfrak m}$ of dimension two. A regular local ring is
a domain. Since $F$ is monic in $x$, the ring $R$ is free over
$V[y]$ and hence flat over $V$. Its localization $A$ is also
$V$-flat, so $V$ embeds into $A$. Together with
$A/\mathfrak m_A=\F_5$, this proves mixed characteristic $(0,5)$.
Expanding $F=0$ gives
\[
 \pi=x^5-5x^4+10x^3-10x^2+5x-xy^4+y^5.
\]
In particular, $\pi\in(x,y)A$ and $\mathfrak m_A=(x,y)A$.
Since $5$ is divisible by $\pi$, the term $5x$ is in
$(x,y)^2A$, as are all remaining terms on the right. Hence
$\pi\in\mathfrak m_A^2$. Since $5=u\pi^4$, also
$5\in\mathfrak m_A^2$, so $A$ is ramified in the usual
mixed-characteristic sense. The reduction of the equation has
no linear term at $m$, so the one-dimensional special fiber is
singular there; $A$ is not smooth over $V$.

Modulo $y$, the equation becomes
\[
 f(t)=(t-1)^5+\zeta
 =t^5-5t^4+10t^3-10t^2+5t-\pi.
\]
This is Eisenstein at $\pi$. The finite free $V$-algebra
$B=V[t]/(f)$ is a domain. Every maximal ideal lies over the
maximal ideal of $V$, and
$B/\pi B=\F_5[t]/(t^5)$, so $B$ is local. Since $f(t)=0$
implies $\pi\in(t)B$, its maximal ideal is $(t)$.
Its dimension is one, making it a DVR. The localization defining
$A/yA$ leaves this local ring unchanged, proving
\eqref{eq:divisor}. Its fraction field has degree five over $L$
and residue degree one, hence total ramification. The constant
term gives
\[
 N_{D/L}(t)=(-1)^5f(0)=\pi.
\]
The relation $(1-t)^5=\zeta$ also follows from $f(t)=0$.
As $S[1/\pi]=D$, the equation $y=0$ on the generic affine
chart gives a single closed point. There is no additional point
with $Y=Z=0$ on $C$, since substitution in $P$ would force
$X=0$. This proves the assertion about $q_0$ on all of $C$.

We have constructed the distinguished point with residue degree one.
It remains to establish the residue-degree constraint on points that
specialize away from $m$.
A closed point of $C$ has residue field finite over $L$.
Properness of $\X/V$ extends its $E$-point to $\Spec\OO_E$.
On the special fiber $Y=0$ forces $X=0$, so $m$ is the only
point with $Y=0$. If the closed image is elsewhere, $X/Y$ in
$k_E$ is a root of $s^5-s+1$. Thus $k_E$ contains $\F_{5^5}$,
and $5\mid[k_E:\F_5]$.
\end{proof}
